\begin{abstract}
Flow Matching (FM) has emerged as the dominant paradigm for state-of-the-art image and video generation, yet its reliance on full attention over lengthy sequences incurs significant latency. While existing caching methods reduce redundant computations, they yield suboptimal performance, fundamentally constrained by inflexible spatial granularity and empirical, model-specific heuristics.

To bridge this gap, we ground the caching mechanism in the geometric properties of the FM velocity field. We theoretically establish that the \textbf{curvature} inherently dictates the coupling between trajectory fluctuation and computational requirements, thereby providing a rigorous metric to quantify caching staleness loss and automatically derive adaptive granularity and optimal update intervals.

Translating this theoretical insight into practical acceleration, we present \cubic, a train-free caching framework. It bypasses expensive curvature computations via a zero-overhead proxy signal that dynamically guides spatial partitioning and budget-constrained update scheduling. At the system level, a sparse forward engine translates these optimal policies into efficient runtime, ensuring that theoretical guarantees directly translate into actual FLOP savings and linear speedup. Being model-agnostic and virtually hyperparameter-free, \cubic consistently outperforms state-of-the-art caching baselines in quality at customizable acceleration ratios (averaging 2.5$\times$) across Wan2.1, Flux.1, and FireRed-Image-Edit.
\end{abstract}